% 替换版本说明段；保留现有公共 setup、图及 subsection 标签。
\noindent\textbf{构建版本与统计方法。}
当前构建系列基于 \texttt{f856c7b}：两种续花模式各运行三轮
100 跳链，另运行三轮代表性的暂停适配处理 P。九轮使用同一节点构建；
新增计时点仅位于 \texttt{payctl}，六轮连续链均已启用，
未修改生产节点路径。\texttt{721800c} 的完整 N/R/C/P 矩阵
保留为历史结果；历史读取实验仍标注 \texttt{721800c}，
注入延迟实验仍标注 \texttt{9879f64}。
连续链报告三轮总时长的中位数与范围；混合负载报告三轮轮内统计量
的中位数。阶段计时先逐笔求时间差，在各轮内分别计算 P50/P95，
再对三轮对应统计量取中位数，不混池计算。

\subsection{无需逐跳结算等待的连续付款}

三个相互独立的 NoSync 钱包在线构造 100 跳付款链，
循环转移 100~CAL；成员采用同步提交，费用使用用户的最终 FUEL。
fast 模式在认证到账后继续，wait 模式还要求驱动观察到钱包已验证的
公共执行，才构造下一跳。两者均测量
$\mathsf{Sent}_1 \rightarrow \mathsf{Ready}_{100}$，
包含后继构造与重试，不包含首笔付款的构造。
wait 模式增加的是 99 次跳间等待，不包含最后一笔的公共确认。

两种模式的链时长中位数分别为 4.707\,s 和 64.479\,s，
中位数之比为 13.70；三轮范围分别为 4.700--4.726\,s
和 64.387--64.670\,s。fast 为 300 笔付款发起 300 次尝试，
wait 发起 414 次。全部 297 个 fast 后继均使用前一笔的 TXCer，
且首次发送早于四个副本中最早的父交易应用提交完成时间，
最小提前量为 183.439\,ms。
三轮实际产生的子交易先于来源执行的义务数分别为 2、4 和 1；
仅提前发送不计作错序执行。

全部 600 笔付款完成闭合，并通过金额、输入、副本一致性、
费用及 outbox 审计。以首次发送为起点，fast 模式第 100 笔付款的
公共观察时间中位数为 5.477\,s，全部付款的成员闭合时间中位数
为 5.526\,s。前者专指第 100 笔，不以全部付款公共观察时间的最大值
替代其定义。每条链支出 9,400~FUEL。
对于该固定交易形状，TXCer 为 779~B；
首跳请求为 1,979~B，后续请求为 2,766~B。

\noindent\textbf{客户端阶段计时。}
记录按接收钱包、区块高度和付款事实匹配。
阶段分解采用取数端 \texttt{committee0} 的应用提交完成时间，
而非上述提前量检查所用的四副本最早时间。
一般阶段每轮包含 100 个付款样本，跳间过渡每轮包含 99 个；
六轮合计分别为 600 和 594 个，每种模式分别为 300 和 297 个。
同一接收钱包、同一高度的多个付款复用观察点，并按付款展开统计，
不将其视为相互独立的区块高度样本。

对于目标高度 $h$，\texttt{BlockData} 先取得 $h+1$ 的签名头，
再获取 $h$ 的区块和执行结果。获取 $h+1$ 提交证据时，
RPC 对最新高度读取 \texttt{LoadSeenCommit}，
对较早高度读取 \texttt{LoadBlockCommit}；
接口不要求额外等待 $h+2$。
因此，从应用提交到完整取数返回仍是包含后继证据、
副本推进、轮询和取数的联合区间，不能解释为纯查询开销，
成功返回也不标识证据最早可得的时刻。

wait 模式下，该联合区间的 P50/P95 为
285.884/314.509\,ms。
另行观察的取数后验证、准备与本地记录、
以及事务完成后回调至驱动观察三个区间，
P95 分别为 0.236、1.343 和 4.713\,ms。
验证后的检查点位于 \texttt{VerifyBlock} 成功返回后的
\texttt{PrepareBlock} 入口；
回调发生在本地事务完成后，不等待耐久刷盘。
单次成功获取调用的耗时与上述联合区间重叠，单独报告而不相加；
各阶段分位数同样不相加。

跨进程区间使用同机 Unix 纳秒时间戳相减，
整链时长使用单调时钟。
驱动轮询、跟随器轮询和 commit 设置仍分别为
5、25 和 250\,ms。
fast 后继可以早于公共观察开始构造，
因此负的“公共观察至下一跳构造”差值表示执行重叠，
不是负等待时间。
13.70 是所测配置下两种端到端模式的时长比，
包含重试和证据观察。

\subsection{代表性混合负载下的来源闭合}

十四服务的 P 处理在 70\,s 内按每秒 100 笔安排
7,000 笔普通付款，同时运行八组相互独立的跨组织父子付款，
费用使用明确授权的组织 FUEL。
实验扣留父交易的公开提交及 \texttt{INSTALL}，
关闭自动来源补交，并暂停全部适配端点；
共识和赔付决定继续运行。
每个迟到来源在观察到其赔付后单独释放。
普通流量覆盖注入故障及恢复的全部时段。

三轮共 21,048 笔付款全部闭合，其中包括
21,000 笔普通付款和 24 组父子付款。
每轮赔付并收回 800~CAL。
恢复适配之前，四个委员会副本均报告八个目标已
\texttt{Recovered}，而对应表示的
\texttt{Committed} 和 \texttt{Materialized} 标志均为 false。
这 96 次副本--目标检查确认，在所测执行中，
赔付和回款均先于适配恢复完成。
每轮委员会状态哈希在共同高度一致，
CAL/FUEL、输入、费用、预留及 outbox 审计全部通过。
每轮费用为 659,544~FUEL，其中
589,384 为奖励，70,160 被销毁。

表~\ref{tab:current-p} 同时报告普通付款到账延迟、
发送前回压和发送后排空。
轮内到账 P50/P95 的三轮中位数分别为
54.538/147.549\,ms；
公共观察和成员观察 P95 的三轮中位数分别为
2.187 和 2.255\,s。
三轮均触及 256 笔未完成付款的驱动上限。
发送滞后 P95 最高为 780.727\,ms，
计划时刻至到账 P95 最高为 851.742\,ms；
后者逐笔计算，并非两个分位数相加。
实际发送跨度为 70.080--70.771\,s，
最后一笔普通付款发送后再经 0.955--1.115\,s 排空。
该代表性 P 系列展示了普通流量期间、适配不可用时的经济闭合，
不构成当前构建下 N/R/C/P 的干扰对照。

\begin{table*}[t]
\centering
\caption{构建 \texttt{f856c7b} 上代表性 P 处理的普通付款结果。
每轮包含 7,000 笔普通付款及八组父子付款。
到账从实际发送起计；发送滞后从计划时刻计至实际发送；
计划至到账直接逐笔计算。排空从最后一笔普通付款的实际发送起计。
末行分别取各列三轮统计量的中位数。}
\label{tab:current-p}
\small
\setlength{\tabcolsep}{5pt}
\begin{tabular}{lrrrrrr}
\toprule
轮次 &
\shortstack{到账 P50\\(ms)} &
\shortstack{到账 P95\\(ms)} &
\shortstack{发送滞后 P95\\(ms)} &
\shortstack{计划至到账 P95\\(ms)} &
\shortstack{发送跨度\\(s)} &
\shortstack{排空\\(s)} \\
\midrule
P-1 & 54.538 & 156.680 &  90.335 & 189.591 & 70.080 & 1.090 \\
P-2 & 54.007 & 147.549 & 245.315 & 326.828 & 70.235 & 0.955 \\
P-3 & 55.001 & 141.424 & 780.727 & 851.742 & 70.771 & 1.115 \\
\midrule
中位数 & 54.538 & 147.549 & 245.315 & 326.828 & 70.235 & 1.090 \\
\bottomrule
\end{tabular}
\end{table*}